\begin{question}{Properties of Finite-state Language Models}

\begin{subquestion}

\subsubsection*{Approach}

We define the duplicated automaton $\mathcal{A}_{dup} = (\Sigma, Q_{dup}, \delta_{dup}, \lambda_{dup}, \rho_{dup})$ as follows:

\begin{description}

    \item[$\Sigma$:]
    The alphabet remains unchanged. The duplicated automaton reads the exact same symbols as the original one.

    \item[$Q_{dup}$:]
    The state space is doubled to create two parallel layers. We define this mathematically as $Q_{dup} = Q \times \{1, 2\}$, where $1$ and $2$ represent the layer numbers.

    \item[$\delta_{dup}$:]
    For every original transition $q \xrightarrow{a, w} q'$ in $\mathcal{A}$, we construct the transitions in $\mathcal{A}_{dup}$ as follows:

    \begin{itemize}
        \item \textbf{Layer 1 Traversal:} A transition $(q, 1) \xrightarrow{a, w} (q', 1)$.
        \item \textbf{Layer 2 Traversal:} A transition $(q, 2) \xrightarrow{a, w} (q', 2)$.
        \item \textbf{The One-Way Bridge:} A transition $(q, 1) \xrightarrow{a, w} (q', 2)$.
    \end{itemize}

    \item[$\lambda_{dup}$:]
    We force all paths to begin in the first layer. We do this by keeping the original initial weights for Layer 1 states and setting the initial weights for all Layer 2 states to $0$.

    \item[$\rho_{dup}$:]
    We force all paths to finish in the second layer. We do this by keeping the original final weights for Layer 2 states and setting the final weights for all Layer 1 states to $0$.

\end{description}

To find the mean string length under $p_{LM}$, we simply execute the \texttt{allsum} algorithm from the notes on $\mathcal{A}_{dup}$.

\subsubsection*{Intuition}

In our two-layer automaton, any path $\pi$ from the original automaton must cross the one-way bridge to Layer 2 exactly once. Because the path has $|\pi|$ steps, there are $|\pi|$ different opportunities to make this jump. 

Therefore, the \texttt{allsum} algorithm discovers $|\pi|$ distinct routes for the original path. By adding up all valid routes, the algorithm naturally adds the path's probability to the total sum $|\pi|$ times, seamlessly computing $\text{Probability} \times \text{Length}$ to find the mean.

\subsubsection*{Runtime}

The duplicated automaton $\mathcal{A}_{dup}$ has twice as many states as the original automaton. If we plug this new size into the runtime for the algorithm from the notes, we get:
$$ \mathcal{O}((2|Q|)^3) = \mathcal{O}(8|Q|^3) $$

We can drop the constant multiplier, because it doesn't change the fundamental growth rate of the algorithm. This leaves us with our final time complexity:
$$ \mathcal{O}(|Q|^3) $$

\end{subquestion}

\begin{subquestion}

\subsubsection*{Approach}

To compute the variance, we use the identity $\text{Var}(X) = E[X^2] - E[X]^2$. Having already determined $E[X]$ via the duplicated automaton, we now construct a three-layered automaton $\mathcal{A}_{tri} = (\Sigma, Q_{tri}, \delta_{tri}, \lambda_{tri}, \rho_{tri})$ to find $E[X^2]$.

The state space is defined as $Q_{tri} = Q \times \{1, 2, 3\}$.

For every original transition $q \xrightarrow{a, w} q'$ in $\mathcal{A}$, the weights in $\mathcal{A}_{tri}$ are assigned as follows:

\subsubsection*{Transitions of $\mathcal{A}_{tri}$}

For every original transition $q \xrightarrow{a, w} q'$ in $\mathcal{A}$, we define the following transitions in $\mathcal{A}_{tri}$:

\begin{itemize}
    \item $(q, i) \xrightarrow{a, w} (q', i)$ for $i \in \{1, 2, 3\}$
    \item $(q, 1) \xrightarrow{a, w} (q', 3)$
    \item $(q, 1) \xrightarrow{a, 2w} (q', 2)$
    \item $(q, 2) \xrightarrow{a, w} (q', 3)$
\end{itemize}

By running \texttt{allsum} on $\mathcal{A}_{tri}$, we obtain $E[X^2]$.

\subsubsection*{Intuition}

The second moment $E[X^2]$ is the expectation of the squared path length, $\sum p(\pi) \cdot |\pi|^2$. By expanding the square as a double summation, $|\pi|^2 = \sum_{i} \sum_{j} 1$, we see that we must account for every possible pair of steps in a path. 

The three layers facilitate this by acting as a "selection" mechanism:
\begin{enumerate}
    \item The path starts in Layer 1.
    \item It moves to Layer 2 when the first step of a pair is "selected."
    \item It moves to Layer 3 when the second step is "selected."
\end{enumerate}
The double-step bridge ensures that steps can be paired with themselves.

\subsubsection*{Weighting Logic}

By using these weights, the \texttt{allsum} result for $\mathcal{A}_{tri}$ yields the second moment $E[X^2]$. Specifically, for a path $\pi$, the weighting ensures that the off-diagonal terms of the expansion $(\sum 1)^2$ are counted twice and the diagonal terms are counted once, successfully yielding $|\pi|^2$.

\subsubsection*{Runtime}

The state space for $\mathcal{A}_{tri}$ is $3|Q|$. The time complexity for the \texttt{allsum} algorithm on this expanded graph is:
$$\mathcal{O}((3|Q|)^3) = \mathcal{O}(27|Q|^3)$$

Dropping the constant coefficient, the final complexity remains:
$$\mathcal{O}(|Q|^3)$$

\end{subquestion}

\begin{subquestion}

To prove that $H(p_{LM}) = \lambda^{\top}(I - T)^{-1}\xi$, we utilize the provided hints regarding the local transition entropy and the chain rule for entropy[cite: 173, 175, 183].

\subsubsection*{1. Defining the Variables}
Following the suggestion in the prompt, let $\Pi_q$ be a random variable ranging over paths in $\mathcal{A}$ starting in state $q$[cite: 179, 180]. Let $\tau_q$ be a random variable over the outgoing transitions from $q$ or termination at $q$[cite: 179, 180]. 
By definition, the local entropy at state $q$, denoted $\xi_q$, is[cite: 176]:
\begin{equation}
    \xi_q = H(\tau_q) = -\sum_{q \xrightarrow{y/w} q'} w \log w - \rho_q \log \rho_q
\end{equation}

\subsubsection*{2. Recursive Entropy Decomposition}
The entropy of a path starting at state $q$ can be decomposed using the chain rule of entropy: $H(X, Y) = H(X) + H(Y|X)$[cite: 183, 184]. For a path $\Pi_q$, the "first step" is $\tau_q$. Since the automaton is deterministic, the next state $q'$ is uniquely determined by the choice of transition in $\tau_q$[cite: 172]. Thus:
\begin{equation}
    H(\Pi_q) = H(\tau_q) + \sum_{q \xrightarrow{y/w} q'} \mathbb{P}(\tau_q = q \xrightarrow{y/w} q') H(\Pi_{q'})
\end{equation}
Substituting $\xi_q$ and the transition weights $w$[cite: 182], we obtain:
\begin{equation}
    H(\Pi_q) = \xi_q + \sum_{q \xrightarrow{y/w} q'} w H(\Pi_{q'})
\end{equation}

\subsubsection*{3. Matrix Formulation and Solution}
Let $\mathbf{h} \in \mathbb{R}^{|Q|}$ be the vector of path entropies where $\mathbf{h}_q = H(\Pi_q)$. Using the transition matrix $T$ defined in Eq. (15) [cite: 141, 142, 145], we can rewrite the recursive relation in matrix form:
\begin{equation}
    \mathbf{h} = \xi + T\mathbf{h}
\end{equation}
Rearranging the terms gives $(I - T)\mathbf{h} = \xi$. Since all eigenvalues of $T$ have magnitude less than 1, the matrix $(I - T)$ is invertible[cite: 146]:
\begin{equation}
    \mathbf{h} = (I - T)^{-1}\xi
\end{equation}

\subsubsection*{4. Conclusion}
The total entropy of the language model $H(p_{LM})$ is the expected entropy of paths starting from the initial state distribution $\lambda$[cite: 140, 147, 152]:
\begin{equation}
    H(p_{LM}) = \sum_{q \in Q} \lambda_q H(\Pi_q) = \lambda^{\top}\mathbf{h}
\end{equation}
Substituting the expression for $\mathbf{h}$ yields the desired result[cite: 175]:
\begin{equation}
    H(p_{LM}) = \lambda^{\top}(I - T)^{-1}\xi
\end{equation}

\end{subquestion}

\end{question}